\documentclass[11pt]{amsart}

\usepackage[letterpaper,margin=1in]{geometry}
\usepackage{amsmath,amssymb,mathtools}
\usepackage{booktabs}
\usepackage{microtype}
\usepackage{xcolor}
\usepackage[colorlinks=true,linkcolor=blue!45!black,citecolor=blue!45!black,
  urlcolor=blue!45!black]{hyperref}

\newcommand{\K}{\mathbb{K}}
\newcommand{\Sym}{\operatorname{Sym}}
\newcommand{\Span}{\operatorname{span}}
\newcommand{\im}{\operatorname{im}}
\newcommand{\rank}{\operatorname{rank}}
\newcommand{\perm}{\operatorname{perm}}
\newcommand{\CR}{\mathrm{R}_{\mathrm{Chow},\K}}
\newcommand{\cD}{\mathcal{D}}

\newtheorem{theorem}{Theorem}[section]
\newtheorem{proposition}[theorem]{Proposition}
\newtheorem{lemma}[theorem]{Lemma}
\newtheorem{corollary}[theorem]{Corollary}
\theoremstyle{definition}
\newtheorem{definition}[theorem]{Definition}
\newtheorem{remark}[theorem]{Remark}

\title[A Candidate Proof for the Six-by-Six Permanent]
{A Candidate Proof for the Ordinary Chow Rank\\
of the Six-by-Six Permanent}
\author{}
\date{August 21, 2026}
\subjclass[2020]{Primary 15A69; Secondary 14N05, 68V05}
\keywords{permanent, Chow rank, product rank, catalecticant,
computer-assisted proof}

\begin{document}

\begin{abstract}
Let \(\perm_6\) be the permanent of a generic \(6\) by \(6\) matrix.
We record a conditional project proof candidate for the assertion that, over
an algebraically closed field of characteristic zero, its ordinary Chow rank
is \(32\).  The upper bound is
Glynn's \(32\)-term identity.  The lower bound combines a symmetric
image-span inequality with a half-defect estimate for quotient derivative
symbols of a degree-six Chow term.  The finite rational rows used in the
half-defect estimate and the final cancellation are recorded explicitly,
together with the exact replay boundary.  No assertion about border Chow
rank or the general \(n\) formula is made.
\end{abstract}

\maketitle

\section{Statement and notation}

Let \(\K\) be an algebraically closed field of characteristic zero and let
\[
 V=\Span_{\K}\{x_{ij}:1\leq i,j\leq6\},
 \qquad \dim V=36.
\]
The six-by-six permanent is
\[
 \perm_6=\sum_{\sigma\in S_6}\prod_{i=1}^{6}x_{i,\sigma(i)}
 \in\Sym^6(V).
\]

\begin{definition}
A degree-six \emph{Chow term} is a product
\[
 T=\ell_1\ell_2\cdots\ell_6,
 \qquad \ell_i\in V.
\]
The ordinary Chow rank \(\CR(f)\) of \(f\in\Sym^6(V)\) is the least
integer \(N\) for which \(f\) is a sum of \(N\) Chow terms.
\end{definition}

For a homogeneous polynomial \(f\), write \(\cD_q(f)\subseteq\Sym^q(V)\)
for the span of all derivatives of \(f\) having degree \(q\).  Equivalently,
\(\cD_q(f)\) is the image of the appropriate catalecticant.  Put
\[
 E_q=\cD_q(\perm_6).
\]
The squarefree subpermanents give
\[
 \dim E_q=\binom{6}{q}^2;
 \qquad \dim E_2=225,\quad \dim E_3=400.
\]
In particular, the middle catalecticant of \(\perm_6\) has rank \(400\).

\begin{theorem}\label{thm:main}
Assume Proposition~\ref{prop:half-defect}.  Then, over \(\K\),
\[
 \boxed{\CR(\perm_6)=32.}
\]
\end{theorem}

The repository does not promote Theorem~\ref{thm:main} to an unconditional
exact-rank result until Proposition~\ref{prop:half-defect} has a complete
independent proof or independently generated exact certificates for every
load-bearing local rank.  The lower bound rests on that proposition.  Its statement
must be read literally: it covers arbitrary factor dependencies, repeated
factors, and arbitrary quotients of the actual factor span.

\section{A symmetric image-span inequality}

\begin{lemma}[Symmetric image-span lemma]\label{lem:image-span}
Let \(A_i:W^*\to W\), \(1\leq i\leq m\), be symmetric linear maps.  Set
\[
 r_i=\rank A_i,
 \qquad
 D=\dim\sum_{i=1}^{m}\im A_i.
\]
Then
\[
 \rank\!\left(\sum_{i=1}^{m}A_i\right)
 \geq 2D-\sum_{i=1}^{m}r_i.
\]
\end{lemma}

\begin{proof}
Choose a surjection \(B_i:\K^{r_i}\to\im A_i\subseteq W\).  Symmetry
implies \(\ker A_i=(\im A_i)^\perp\), so the form induced by \(A_i\) on
its rank quotient is nondegenerate.  Consequently
\[
 A_i=B_iJ_iB_i^*
\]
for an invertible symmetric \(r_i\) by \(r_i\) matrix \(J_i\).  Let
\[
 B=[B_1\ \cdots\ B_m],
 \qquad J=\operatorname{diag}(J_1,\ldots,J_m),
 \qquad R=\sum_i r_i.
\]
Then \(\sum_iA_i=BJB^*\), while \(\rank B=\rank B^*=D\) and
\(\rank(BJ)=D\).  Sylvester's rank inequality gives
\[
 \rank(BJB^*)\geq\rank(BJ)+\rank(B^*)-R=2D-R.
\]
\end{proof}

\section{The required local half-defect estimate}

Let \(T=\ell_1\cdots\ell_6\neq0\) be a Chow term and set
\[
 L=\Span\{\ell_1,\ldots,\ell_6\},\qquad
 U=\cD_3(T),\qquad u=\dim U,
\]
\[
 F=\cD_2(T),\qquad R=F\cap E_2.
\]
Given a quotient \(P:L\twoheadrightarrow D\) of rank \(d\), differentiating
cubics in \(U\) and then reducing their quadratic derivatives modulo \(R\)
defines the quotient derivative symbol
\[
 \beta_{P,R}:U\longrightarrow D\otimes(F/R).
\]
Only the rank of this natural map is used below.

\begin{proposition}[Required half-defect quotient-symbol estimate]
\label{prop:half-defect}
For every nonzero degree-six Chow term \(T\) and every quotient
\(P:L\twoheadrightarrow D\) of rank \(d\),
\begin{equation}\label{eq:HD}
 \rank\beta_{P,R}+\frac{20-u}{2}\geq\frac{10}{3}d.
\end{equation}
\end{proposition}

\begin{remark}[Current reduction of the proof obligation]
The candidate argument divides the proof by the factor-span dimension
\(\ell=\dim L\).

For \(\ell\leq4\), a diagonal one-parameter subgroup degenerates \(T\) to
a monomial in \(\ell\) active variables with all exponents positive.  The
corresponding middle-catalectic ranks have floors
\[
 1,\ 2,\ 4,\ 8
\]
for \(\ell=1,2,3,4\), respectively.  The intersection of the resulting
quadratic derivative space with \(E_2\) has dimension at most one.  Combining
this with the kernel-preimage estimate for intermediate quotient ranks and
the full-quotient injectivity criterion gives the following adjusted rows:
\[
\begin{array}{c|c}
\toprule
\ell & \rank\beta_{P,R}+(20-u)/2\text{ for }d=0,\ldots,\ell\\
\midrule
1&(19/2,\ 21/2)\\
2&(8,\ 9,\ 11)\\
3&(5,\ 15/2,\ 9,\ 12)\\
4&(0,\ 9/2,\ 8,\ 10,\ 14)\\
\bottomrule
\end{array}
\]
The full-quotient assertion uses \(E_2^{(1)}=E_3\): a kernel vector would
produce a nonzero permanent cubic supported on fewer than nine essential
variables, which is impossible.

For \(\ell=6\), the proposed quotient-symbol reduction gives
\begin{equation}\label{eq:ell6-row}
 (0,\ 7,\ 10,\ 10,\ 15,\ 17,\ 20),
 \qquad d=0,\ldots,6.
\end{equation}
For \(\ell=5\), linear changes of basis reduce the repeated-factor patterns
to
\[
 x_1x_2x_3x_4x_5(x_1+\cdots+x_s),\qquad1\leq s\leq5.
\]
Their proposed adjusted symbol rows are
\[
\begin{array}{c|c}
\toprule
s & \rank\beta_{P,R}+(20-u)/2\text{ for }d=0,\ldots,5\\
\midrule
5&(0,\ 7,\ 10,\ 10,\ 15,\ 17)\\
4&(0,\ 8,\ 12,\ 13,\ 16,\ 19)\\
3&(1,\ 8,\ 11,\ 11,\ 16,\ 18)\\
1,2&(3,\ 9,\ 9,\ 10,\ 16,\ 17)\\
\bottomrule
\end{array}
\]
The last row is intended to use the seven-dimensional directional-shadow floor and the
one-dimensional kernel gap at quotient rank four.  Direct comparison of
every entry in all three tables with \(10d/3\) proves \eqref{eq:HD}
provided the geometric reductions and the displayed exact symbol rows are
established independently.  Those derivations are the unresolved proof
obligation in this candidate package.
\end{remark}

\begin{remark}[Exact replay boundary]\label{rem:boundary}
The project replay script checks every rational comparison in the displayed
rows and the final cancellation in Section~\ref{sec:global}.  It does not
derive the geometric degeneration, the normal-form reduction, the
kernel-preimage estimate, or the exact quotient-symbol ranks
\eqref{eq:ell6-row}.  Those are mathematical inputs to the replay, not
consequences of the JSON certificate.
\end{remark}

\section{The global lower bound}\label{sec:global}

Assume that
\begin{equation}\label{eq:decomp}
 \perm_6=\sum_{i=1}^{N}T_i
\end{equation}
with all \(T_i\neq0\).  Let
\[
 A_i=C_{3,3}(T_i):\Sym^3(V^*)\longrightarrow\Sym^3(V)
\]
be the symmetric middle catalecticant and put
\[
 U_i=\im A_i,\qquad u_i=\dim U_i,
 \qquad \delta_i=20-u_i,\qquad \Delta=\sum_i\delta_i.
\]
Set \(U=\sum_iU_i\).  Since \(E_3\subseteq U\), write
\[
 h=\dim(U/E_3),\qquad \dim U=400+h.
\]

The sum \(\sum_iA_i\) is the middle catalecticant of \(\perm_6\), hence
has rank \(400\).  Lemma~\ref{lem:image-span} yields
\[
 400\geq2(400+h)-\sum_i u_i
 =800+2h-(20N-\Delta).
\]
Therefore
\begin{equation}\label{eq:h-upper}
 \boxed{h\leq10N-200-\frac{\Delta}{2}.}
\end{equation}

We next derive the opposite estimate.  Let
\[
 F=\sum_i\cD_2(T_i),\qquad Q=F/E_2.
\]
Differentiating \eqref{eq:decomp} twice gives \(E_2\subseteq F\).  The
blockwise quotient derivative symbol is
\[
 \widetilde\beta:\bigoplus_iU_i\longrightarrow V\otimes Q.
\]
The prolongation identity \(E_2^{(1)}=E_3\) implies
\[
 \ker\widetilde\beta
 =\left\{(g_i):\sum_i g_i\in E_3\right\}.
\]
Thus \(\im\widetilde\beta\) is naturally isomorphic to \(U/E_3\), and
\begin{equation}\label{eq:beta-rank}
 \rank\widetilde\beta=h.
\end{equation}
This identity already accounts for all overlaps among the spaces \(U_i\).

Let \(L_i\) be the factor span of \(T_i\).  Since \(E_2\subseteq F\) and
the first derivatives of \(E_2\) span \(V\), the spaces \(L_i\) together
span all \(36\) variables.  Choose an arbitrary ordering and define
\[
 W_i=\sum_{j\leq i}L_j,
 \qquad d_i=\dim(W_i/W_{i-1}).
\]
Then
\begin{equation}\label{eq:sum-di}
 \sum_i d_i=36.
\end{equation}

Project the first tensor factor of \(V\otimes Q\) to
\((W_i/W_{i-1})\otimes Q\).  This kills every earlier symbol block.  On
the current block it is induced by the arbitrary quotient
\[
 P_i:L_i\longrightarrow L_i/(L_i\cap W_{i-1}),
 \qquad \rank P_i=d_i.
\]
Moreover, the natural map
\[
 \cD_2(T_i)/(\cD_2(T_i)\cap E_2)\longrightarrow F/E_2=Q
\]
is injective.  Hence the projected local ranks add along this filtration.
Proposition~\ref{prop:half-defect}, \eqref{eq:beta-rank}, and
\eqref{eq:sum-di} give
\begin{align}
 h
 &\geq\sum_i\left(\frac{10}{3}d_i-\frac{\delta_i}{2}\right)\\
 &=\frac{10}{3}\cdot36-\frac{\Delta}{2}.
\end{align}
Therefore
\begin{equation}\label{eq:h-lower}
 \boxed{h\geq120-\frac{\Delta}{2}.}
\end{equation}

Comparing \eqref{eq:h-upper} and \eqref{eq:h-lower}, the complete
individual middle-rank defect cancels:
\[
 120-\frac{\Delta}{2}
 \leq10N-200-\frac{\Delta}{2}.
\]
It follows that
\begin{equation}\label{eq:lower32}
 N\geq32.
\end{equation}

\section{Glynn's upper bound and conclusion}

For a generic \(n\) by \(n\) matrix \(X=(x_{ij})\), Glynn's identity
\cite{Glynn2010} is
\begin{equation}\label{eq:glynn}
 \perm_n
 =2^{1-n}
 \sum_{\substack{\varepsilon\in\{\pm1\}^n\\\varepsilon_1=1}}
 \left(\prod_{i=1}^{n}\varepsilon_i\right)
 \prod_{j=1}^{n}\left(\sum_{i=1}^{n}\varepsilon_i x_{ij}\right).
\end{equation}
There are \(2^{n-1}\) summands, and every summand is a product of \(n\)
linear forms.  For \(n=6\), \eqref{eq:glynn} gives
\[
 \CR(\perm_6)\leq32.
\]
Together with \eqref{eq:lower32}, this proves the conditional
Theorem~\ref{thm:main}. \qed

\section{Computer-assisted component}

The finite replay uses exact rational arithmetic only.  In the project root,
the command is
\begin{verbatim}
python scripts/n6_exact_ordinary_chow_rank_32_candidate.py \
  --verify-json data/n6_exact_ordinary_chow_rank_32_candidate.json
\end{verbatim}
The line breaks above are typographical; the script path and JSON path are
each single arguments.  A successful run prints
\begin{verbatim}
PASS: conditional rank-32 arithmetic payload matches
\end{verbatim}
and verifies:
\begin{enumerate}
 \item every displayed half-defect row dominates \(10d/3\);
 \item the lower symbol constant is \(120\);
 \item the upper expression is \(10N-200-\Delta/2\);
 \item \(N=31\) leaves a gap of \(10\), while \(N=32\) has gap zero.
\end{enumerate}
As stated in Remark~\ref{rem:boundary}, this replay certifies the finite
arithmetic interface.  The unresolved geometric and local-rank arguments
remain proof obligations.

\begin{remark}[Scope]
The candidate concerns ordinary Chow rank over characteristic zero.  It does
not determine border Chow rank.  It also does not establish
\(\CR(\perm_n)=2^{n-1}\) for arbitrary \(n\).  Product rank and Chow rank
are used synonymously in the relevant literature; see, for example,
\cite{IltenTeitler2016} for the standard product-rank terminology.
\end{remark}

\begin{thebibliography}{9}

\bibitem{Glynn2010}
D.~G. Glynn,
\emph{The permanent of a square matrix},
European J. Combin. \textbf{31} (2010), no.~7, 1887--1891,
\href{https://doi.org/10.1016/j.ejc.2010.01.010}
{doi:10.1016/j.ejc.2010.01.010}.

\bibitem{IltenTeitler2016}
N.~Ilten and Z.~Teitler,
\emph{Product ranks of the \(3\times3\) determinant and permanent},
Canad. Math. Bull. \textbf{59} (2016), no.~2, 311--319,
\href{https://doi.org/10.4153/CMB-2015-076-1}
{doi:10.4153/CMB-2015-076-1}.

\end{thebibliography}

\end{document}
